\documentclass[14pt,dvipdfmx]{jsarticle}
\input{lecture-preamble.tex}

% 例8の曲線 y=x^3-3x を, 極大・極小の模式図にする. 山が x=a, 谷が x=b
\newcommand{\extremaschema}{%
  \begin{tikzpicture}[x=1.22cm, y=0.83cm, line cap=round, line join=round]
    \draw[-{Stealth}] (-2.30,0) -- (2.30,0) node[right] {$x$};
    \draw[-{Stealth}] (0,-2.45) -- (0,2.50) node[above] {$y$};
    \node[inner sep=1pt] at (0.28,0.32) {O};
    \draw[thick, domain=-2:2, samples=80] plot (\x, {\x*\x*\x-3*\x});
    \draw[dashed, gray] (-1,2) -- (-1,0);
    \draw[dashed, gray] (1,-2) -- (1,0);
    \fill (-1,2) circle (1.6pt);
    \fill (1,-2) circle (1.6pt);
    \node[below=1pt, inner sep=1pt] at (-1,0) {$-1$};
    \node[below=1pt, inner sep=1pt] at (1,0) {$1$};
  \end{tikzpicture}%
}

% 例題4: y=x^3-3x^2+3. 極大 (0,3), 極小 (2,-1)
\newcommand{\reidafourfig}{%
  \begin{tikzpicture}[x=0.96cm, y=0.60cm, line cap=round, line join=round]
    \draw[-{Stealth}] (-1.35,0) -- (3.45,0) node[right] {$x$};
    \draw[-{Stealth}] (0,-2.05) -- (0,4.35) node[above] {$y$};
    \node[inner sep=1pt] at (0.30,-0.40) {O};
    \draw[thick, domain=-1.05:3.12, samples=80] plot (\x, {\x*\x*\x-3*\x*\x+3});
    \draw[dashed, gray] (2,-1) -- (2,0);
    \draw[dashed, gray] (2,-1) -- (0,-1);
    \fill (0,3) circle (1.5pt);
    \fill (2,-1) circle (1.5pt);
    \node[below=1pt, inner sep=1pt] at (2,1) {$2$};
    % 山は y 軸上なので, 3 は曲線の上側に置く. -1 は破線の下
    \node[anchor=south east, inner sep=1pt] at (-0.12,3.12) {$3$};
    \node[anchor=north west, inner sep=1pt] at (-0.8,-1.12) {$-1$};
  \end{tikzpicture}%
}

% 例9: y=x^3. 原点で接線は水平だが, 増減は変わらない
\newcommand{\purecubicfig}{%
  \begin{tikzpicture}[x=1.08cm, y=0.74cm, line cap=round, line join=round]
    \draw[-{Stealth}] (-1.65,0) -- (1.85,0) node[right] {$x$};
    \draw[-{Stealth}] (0,-1.75) -- (0,2.05) node[above] {$y$};
    \node[inner sep=1pt] at (0.48,0.58) {O};
    \draw[thick, domain=-1.22:1.22, samples=60] plot (\x, {\x*\x*\x});
    \draw[dashed, gray] (1,1) -- (1,0);
    \draw[dashed, gray] (1,1) -- (0,1);
    \fill (0,0) circle (1.4pt);
    \fill (1,1) circle (1.4pt);
    \node[below, inner sep=1pt] at (1,0) {$1$};
    \node[anchor=east, inner sep=1pt] at (-0.08,1) {$1$};
  \end{tikzpicture}%
}

% 例10: y=-x^3-x. 常に減少
\newcommand{\deccubicfig}{%
  \begin{tikzpicture}[x=1.08cm, y=0.56cm, line cap=round, line join=round]
    \draw[-{Stealth}] (-1.55,0) -- (1.7,0) node[right] {$x$};
    \draw[-{Stealth}] (0,-2.85) -- (0,2.85) node[above] {$y$};
    \node[inner sep=1pt] at (0.50,0.58) {O};
    \draw[thick, domain=-1.15:1.15, samples=60] plot (\x, {-\x*\x*\x-\x});
    \draw[dashed, gray] (1,-2) -- (1,0);
    \draw[dashed, gray] (1,-2) -- (0,-2);
    \fill (1,-2) circle (1.4pt);
    \node[anchor=north west, inner sep=1pt] at (1.08,-0.12) {$1$};
    \node[anchor=east, inner sep=1pt] at (-0.08,-2) {$-2$};
  \end{tikzpicture}%
}

% 増減表. 極値のマスには値だけを入れる
% 2-4-A と同じく, 列幅は数字の桁によらず一定
\newif\ifzgblank
\newcommand{\zgcolwidth}{3.4zw}
\newcommand{\zgbox}[1]{\makebox[\zgcolwidth][c]{#1}}
\newcommand{\zgval}[1]{\ifzgblank\fitblankbox{#1}\else $#1$\fi}
\newcommand{\zg}[1]{\zgbox{\zgval{#1}}}
\newcommand{\extprimehead}{$f'(x)$}
\newcommand{\extfunchead}{$f(x)$}
\newcommand{\zgsetup}{%
  \renewcommand{\arraystretch}{1.18}%
  \renewcommand{\fitblankpad}{0.15zw}%
  \setlength{\tabcolsep}{0.28zw}%
  \newcommand{\zgdots}{\zgbox{$\cdots$}}%
}
% 3次で先頭の係数が正: 増加 → 極大値 → 減少 → 極小値 → 増加
% 引数: x1, x2, 極大値, 極小値
\newcommand{\maxminchartbody}[4]{%
  \zgsetup
  \begin{tabular}{|c||c|c|c|c|c|}
    \hline
    $x$ & \zgdots & \zg{#1} & \zgdots & \zg{#2} & \zgdots \\
    \hline
    \extprimehead & \zg{+} & \zg{0} & \zg{-} & \zg{0} & \zg{+} \\
    \hline
    \extfunchead & \zg{\nearrow} & \zg{#3} & \zg{\searrow} & \zg{#4} & \zg{\nearrow} \\
    \hline
  \end{tabular}%
}
\newcommand{\maxminchart}[4]{%
  \begingroup\zgblanktrue\maxminchartbody{#1}{#2}{#3}{#4}\endgroup
}
% 臨界点が1つで, 前後の符号が同じ（極値なし）
% 引数: x0, 左符号, 右符号, 左矢印, 値, 右矢印
\newcommand{\flatincdecbody}[6]{%
  \zgsetup
  \begin{tabular}{|c||c|c|c|}
    \hline
    $x$ & \zgdots & \zg{#1} & \zgdots \\
    \hline
    \extprimehead & \zg{#2} & \zg{0} & \zg{#3} \\
    \hline
    \extfunchead & \zg{#4} & \zg{#5} & \zg{#6} \\
    \hline
  \end{tabular}%
}
\newcommand{\flatincdec}[6]{%
  \begingroup\zgblanktrue\flatincdecbody{#1}{#2}{#3}{#4}{#5}{#6}\endgroup
}

\begin{document}

% 2-4-A の定義・定理・Technique 6.2.1 の続き
\lectureheader[8]{数学Ⅱ}{6}{微分法と積分法}{2}{関数の値の変化}{4-B}%
  {関数の極大・極小}
  {193--195}{増減の変わり目から極値を読み取り, グラフの概形をつかもう}
\setcounter{definition}{1}
\setcounter{theorem}{1}
\setcounter{technique}{1}

% ============ 1ページ目（表・左） ============
\vspace{-2.0zw}%
\begin{reviewbox}{関数の増減}
  \begin{enumerate}[label=\textbf{\arabic*}, leftmargin=1.5zw, itemsep=0.2em, topsep=0.1em]
    \item $f'(x)>\fitblankbf{0}$ となる範囲では増加し,
          $f'(x)<\fitblankbf{0}$ となる範囲では減少する.
    \item 増減は\fitblankbf{\text{増減表}}にまとめる.\\
          $f'(x)=0$ となる $x$ も, その範囲に含まれる.
  \end{enumerate}
\end{reviewbox}

\vspace{0.05em}
\topic{B}{関数の極大・極小}

\noindent
\begin{minipage}[c]{0.44\linewidth}
  例8の $f(x)=x^{3}-3x$ では,\\
  $x=-1$ で増加から減少へ,\\
  $x=1$ で減少から増加へ移る.\\[0.45em]
  一般に, 次のようにいう.
\end{minipage}\hfill
\begin{minipage}[c]{0.54\linewidth}
  \centering\extremaschema
\end{minipage}

\vspace{0.05em}
\begin{definitionbox}{極大・極小}
  {\renewcommand{\fitblankpad}{0.8zw}%
  $x=a$ を境目として増加から減少に移るとき,\\
  $f(x)$ は $x=a$ で\fitblankbf{\text{極大}}であるといい,
  $f(a)$ を\fitblankbf{\text{極大値}}という.\\[0.45em]
  $x=b$ を境目として減少から増加に移るとき,\\
  $f(x)$ は $x=b$ で\fitblankbf{\text{極小}}であるといい,
  $f(b)$ を\fitblankbf{\text{極小値}}という.\\[0.45em]
  極大値と極小値をまとめて\fitblankbf{\text{極値}}という.}
\end{definitionbox}

\vfill
\nextpage

% ============ 2ページ目（表・右） ============
\vspace{-1.4zw}%
\begin{exampleprob}{4}
  関数 $y=x^{3}-3x^{2}+3$ の増減を調べ, 極値を求めよ.
  また, そのグラフをかけ.
\end{exampleprob}
\vspace{0.5em}

\noindent
\begin{minipage}[c]{0.56\linewidth}
  \begin{answerbox}{}[16.5]
    $y'=3x^{2}-6x=3x(x-2)$\\[0.12em]
    $y'=0$ とすると \quad $x=0,\ 2$\\[0.1em]
    $x = 0$ のとき $y = 3$,\quad $x = 2$ のとき $y = -1$
  \end{answerbox}
\end{minipage}\hfill
\begin{minipage}[c]{0.42\linewidth}
  \centering\reidafourfig
\end{minipage}

\vspace{1.5em}
よって, $y$ の増減表は次のようになる.
\par\vspace{1em}
{\centering
  \renewcommand{\extprimehead}{$y'$}%
  \renewcommand{\extfunchead}{$y$}%
  \maxminchart{0}{2}{3}{-1}\par}

\vspace{1.5em}
{\renewcommand{\fitblankpad}{0.7zw}%
したがって,

\qquad \qquad $x=\fitblankbf{0}$ で極大値 $\fitblankbf{3}$ をとり,\quad $x=\fitblankbf{2}$ で極小値 $\fitblankbf{-1}$ をとる.}

\vfill
\begin{techniquebox}{極値の求め方}
  \begin{enumerate}[label=\textbf{\arabic*}, leftmargin=1.5zw, itemsep=0.22em, topsep=0.12em]
    \item 導関数 $\fitblankbf{f'(x)}$ を求める.
    \item $\fitblankbf{f'(x)=0}$ となる $x$ で範囲を分け, 各範囲の符号を調べる.
    \item 増減表で, 増加から減少に変わる $x$ が\fitblankbf{\text{極大}},
          減少から増加に変わる $x$ が\fitblankbf{\text{極小}}である.
  \end{enumerate}
\end{techniquebox}

\nextpage

% ============ 3ページ目（裏・左） ============
\vspace{-1.4zw}%
関数が常に増加, または常に減少する範囲では, 増減が入れかわることはない.
したがって, 関数は\fitblankbf{\text{極値をもたない}}.

\vspace{0.45em}
\noindent{\bfseries 例\,9}\quad
関数 $f(x)=x^{3}$ の増減と極値を調べる.
\par\vspace{0.3em}
\noindent
\begin{minipage}[c]{0.56\linewidth}
  $f'(x)=\fitblankbf{3x^{2}}$\\[0.35em]
  $f'(x)=0$ とすると \quad $x=\fitblankbf{0}$\\[0.45em]
  増減表は次のようになる.
  \par\vspace{0.35em}
  {\centering\flatincdec{0}{+}{+}{\nearrow}{0}{\nearrow}\par}
  \vspace{0.45em}
  $x=0$ の前後で符号は変わらない.
  \par\vspace{0.2em}
  よって, $f(x)$ は常に\fitblankbf{\text{増加}}し,
  極値を\fitblankbf{\text{もたない}}.
\end{minipage}\hfill
\begin{minipage}[c]{0.40\linewidth}
  \centering\purecubicfig
\end{minipage}

\vspace{0.85em}
\noindent{\bfseries 例\,10}\quad
関数 $f(x)=-x^{3}-x$ の増減と極値を調べる.
\par\vspace{0.25em}
\noindent
\begin{minipage}[c]{0.58\linewidth}
  $f'(x)=\fitblankbf{-3x^{2}-1}$\\[0.4em]
  $-3x^{2}\leqq 0$ であるから, 常に $f'(x)\fitblankbf{<0}$.
  \par\vspace{0.35em}
  よって, $f(x)$ は常に\fitblankbf{\text{減少}}し,
  極値を\fitblankbf{\text{もたない}}.
\end{minipage}\hfill
\begin{minipage}[c]{0.40\linewidth}
  \centering\deccubicfig
\end{minipage}

\vspace{0.7em}
\begin{theorembox}{極値と $f'(a)$}
  $x$ の多項式で表される関数 $f(x)$ について,
  \par\vspace{0.2em}
  \hspace*{1zw}$f(x)$ が $x=a$ で極値をとるならば,
  $f'(a)=\fitblankbf{0}$ である.
\end{theorembox}

\vspace{0.15em}
\noindent
\setbox0=\hbox{\textbf{注意}\quad}%
\dimen0=\wd0
\box0\begin{minipage}[t]{\dimexpr\linewidth-\dimen0\relax}
逆は成り立たない. $f'(a)=0$ であっても, $x=a$ で極値をとるとは限らない.\\[0.25em]
例9では $f'(0)=0$ であるが, $f(x)=x^{3}$ は $x=0$ で極値をとらない.
\end{minipage}

\vfill
\nextpage

% ============ 4ページ目（裏・右） ============
\vspace{-1.4zw}%
\noindent{\bfseries 応用例題\,2}\quad
関数 $f(x)=x^{3}+ax+b$ が $x=2$ で極小値 $-6$ をとるように,
定数 $a$, $b$ の値を定めよ. また, 極大値を求めよ.
\par\vspace{0.4em}
\textbf{考え方}\quad
定理 6.2.2 より $f'(2)=\fitblankbf{0}$ であり, さらに $f(2)=-6$ である.
\par\vspace{0.15em}
\hspace*{4.5zw}$f'(2)=0$ だけでは極値をとるとは限らないので, 増減表で確認する.

\vspace{0.15em}
\begin{answerbox}{}[13]
  $f'(x)=3x^{2}+a$\\[0.15em]
  $f'(2)=0$,\ $f(2)=-6$ より \quad $12+a=0$,\ $8+2a+b=-6$\\[0.15em]
  よって \quad $a=-12$,\quad $b=10$\\[0.2em]
  このとき \quad $f(x)=x^{3}-12x+10$\\[0.08em]
  $f'(x)=3x^{2}-12=3(x+2)(x-2)$
  \boxanswer{$a=-12$,\ $b=10$}
\end{answerbox}

\vspace{0.35em}
増減表は次のようになる.
\par\vspace{0.35em}
{\centering
  \maxminchart{-2}{2}{26}{-6}\par}

\vspace{0.45em}
{\renewcommand{\fitblankpad}{0.7zw}%
したがって, $x=\fitblankbf{-2}$ で極大値 $\fitblankbf{26}$ をとり,\\
$x=\fitblankbf{2}$ で極小値 $\fitblankbf{-6}$ をとる.}

\vfill
\begin{summarybox}
  \begin{enumerate}[label=\textbf{\arabic*}, leftmargin=1.5zw, itemsep=0.16em, topsep=0.08em]
    \item 増加から減少に変わるところで\fitblankbf{\text{極大}},
          減少から増加に変わるところで\fitblankbf{\text{極小}}となる.
    \item 極値は\fitblankbf{\text{増減表}}から読み取る.\\
          増減が変わらなければ\fitblankbf{\text{極値はない}}.
    \item 極値をとるならば $f'(a)=\fitblankbf{0}$ である.
          逆は成り立たないので, 増減表で確認する.
  \end{enumerate}
\end{summarybox}

\vspace{0.02em}
\nexttime{極値と定義域の端に注目して, 関数の最大・最小を考えてみよう!}

\end{document}
